\documentclass[11pt]{article}

\usepackage[margin=1.2in]{geometry}
\usepackage[T1]{fontenc}
\usepackage{lmodern}
\usepackage{amsmath,amssymb,amsthm}
\usepackage{booktabs}
\usepackage{microtype}
\usepackage[colorlinks=true,linkcolor=blue!50!black,citecolor=blue!50!black,urlcolor=blue!50!black]{hyperref}
\usepackage{xcolor}

\newtheorem{theorem}{Theorem}[section]
\newtheorem{lemma}[theorem]{Lemma}
\newtheorem{proposition}[theorem]{Proposition}
\newtheorem{corollary}[theorem]{Corollary}
\theoremstyle{definition}
\newtheorem{definition}[theorem]{Definition}
\newtheorem{example}[theorem]{Example}
\theoremstyle{remark}
\newtheorem{remark}[theorem]{Remark}

\newcommand{\lie}[3]{\langle #1,#2,#3\rangle}
\newcommand{\ms}[1]{\{#1\}}
\newcommand{\reach}{\longrightarrow^{*}}
\newcommand{\move}{\longrightarrow}
\newcommand{\N}{\mathbb{N}}
\newcommand{\Z}{\mathbb{Z}}
\newcommand{\Scat}{\mathsf{Scat}}

\title{One False Sum:\\
\large A complete, formally verified characterization of solvability\\
in a splitting--merging game with a single arithmetic lie}
\author{}
\date{July 2026}

\begin{document}

\maketitle

\begin{abstract}
We study a one-player game on finite multisets of positive integers in which a
ball of value $n$ may be split into $\lfloor n/2\rfloor$ and
$\lceil n/2\rceil$, and two balls may be merged into their sum---except that
the arithmetic is governed by a single \emph{false sum} $a+b=c$: the ball $c$
splits only into $a$ and $b$, and the pair $\{a,b\}$ merges only into $c$.
Writing $g=|a+b-c|$, $H=\max(a+b,c)$ and $M=H+1$, we prove that for
\emph{every} false sum $(a,b,c)$ with $a,b,c\ge 1$ and all
$s,t\ge M$, the single ball $s$ can be transformed into the single ball $t$
if and only if $g \mid t-s$; moreover the threshold $M=H+1$ cannot be lowered
to $H$ in general.  The sufficiency proof is a taxonomy: a generic
``all-ones hub'' handles most parameters, but it provably breaks down on
several degenerate families---the \emph{traps}, in which some ball values can
never be reduced to units---each of which requires its own construction,
including one eleven-move gadget for the single most pathological shape
$\lie{c}{2c}{c}$.  A dichotomy lemma (every value $v\ge c$ is either of the
trap form $2^k c$ or ``scatterable'') glues the taxonomy into a complete case
analysis.  All results, from the move semantics up to the final
characterization, are formally verified in Lean~4 (core library only), in
roughly $7{,}500$ lines comprising $328$ theorems with no unproven
assumptions.
\end{abstract}

\section{Introduction}

Imagine a puzzle in which you are handed a single ball labelled with a
positive integer and asked to produce a single ball with a different label.
Two operations are available: \emph{splitting} a ball into two balls whose
labels are the (integer) halves of the original, and \emph{merging} two balls
into one whose label is their sum.  Splitting and merging are exact inverses
of one another in spirit, and with true arithmetic the game is
uninteresting: the total is invariant, so a single ball $s$ can become a
single ball $t$ only when $s=t$.

The game becomes interesting when the arithmetic is wrong.  Fix a single
\emph{false sum}: an assertion $a+b=c$ where possibly $a+b\ne c$.  The player
is committed to this belief.\footnote{The game we formalize, \emph{Ya
Stupid}, takes its name from a 2014 viral video in which a child, asked
``what's nine plus ten?'', confidently answers ``twenty-one.''  The false sum
$9+10=21$ is accordingly called the \emph{Classic} configuration below.}
Consequently the ball $c$ no longer splits into its halves---it splits into
$a$ and $b$---and the pair $\{a,b\}$ no longer merges to $a+b$ but to $c$.
Every false move changes the total by $\pm(a+b-c)$, so the total is no longer
invariant and genuine transformations $s \rightsquigarrow t$ with $s \neq t$
become possible.  Two natural questions arise:

\begin{enumerate}
\item[(Q1)] Exactly which pairs $(s,t)$ are solvable?
\item[(Q2)] Is there a threshold above which the answer depends only on the
obvious congruence obstruction?
\end{enumerate}

Our main result answers both completely, for \emph{every} false sum.

\begin{theorem}[Main; formally verified]\label{thm:main-intro}
Let $a,b,c\ge 1$, let $g=|a+b-c|$, $H=\max(a+b,c)$ and $M=H+1$.  For all
$s,t\ge M$,
\[
\ms{s} \reach \ms{t}
\quad\Longleftrightarrow\quad
g \mid t-s .
\]
(When $g=0$, i.e.\ the ``false'' sum is true, the right-hand side reads
$t=s$.)  Moreover, the threshold $M=H+1$ is optimal as a function of $H$: for
the Classic configuration $9+10=21$ one has $H=21$, the pair $(21,23)$
satisfies the congruence, and yet $\ms{21}\not\reach\ms{23}$.
\end{theorem}

The forward implication is an invariant argument and holds far more
generally (for any finite set of simultaneous false sums, with $g$ their
gcd); the content of the theorem is the converse, \emph{sufficiency}.

What makes sufficiency interesting is that the natural proof strategy is
\emph{false}.  The natural strategy---call it the \emph{all-ones hub}---is to
scatter the ball $s$ into $s$ unit balls, adjust the size of the pile of
units by firing false moves, and then rebuild the target $t$.  For many
parameters this works, and \S\ref{sec:hub} makes it precise.  But the game
contains genuine \emph{traps}: configurations in which certain balls can
never be reduced to units at all.  The simplest is $2+2=2$, where a lone
ball $2$ only ever produces multisets of $2$'s (its only split is the false
one, $2\to\{2,2\}$, and two $2$'s only merge falsely back to a $2$).  More
elaborately, when both legs $a,b$ are of the form $c\cdot 2^k$ (as in
$6+6=3$ or $3+6=3$), every route from a large ball down towards units jams on
\emph{locked} copies of $c$, which refuse to halve.  Each trap family needs
its own escape:

\begin{itemize}
\item for the traps with $c\ge 3$, a hub built on \emph{copies of $c$}
rather than units, driven by a peeling recursion
(\S\ref{sec:traps});
\item for the single most pathological shape $\lie{c}{2c}{c}$
(``$c+2c=c$''), where even the peeling recursion jams at one specific ball
value, a bespoke eleven-move excursion that temporarily doubles the total
(\S\ref{sec:peel4c});
\item for $c=2$, a hub built on copies of $2$, with odd values manufactured
by \emph{splitting} even balls rather than assembling units
(\S\ref{sec:smallc});
\item for the doubly degenerate family $1+c=c$, in which the na\"ive
peel-off-a-unit move is unavailable in principle, a
lose-exactly-one recursion (\S\ref{sec:onecc});
\item for $c=1$, an unobstructed hub (\S\ref{sec:smallc}).
\end{itemize}

A short \emph{dichotomy lemma} (Lemma~\ref{lem:dichotomy}: every $v\ge c$ is
either exactly $2^kc$, or ``scatterable'') proves that this taxonomy is
exhaustive, and a final assembly theorem dispatches every $(a,b,c)$ to its
family.

\paragraph{Formal verification.}
Every result in this paper---the move semantics, the invariant, the
sharpness witness, each family construction, the dichotomy, and the final
characterization---has been formally verified in the Lean~4 proof assistant
\cite{lean4}, using only Lean's core library (no \textsf{Mathlib}).  The
development is a single file of $7{,}462$ lines containing $328$ theorems;
it compiles in under a minute and depends only on the three standard
axioms of Lean's classical logic (propositional extensionality, choice,
quotient soundness)---in particular there are no unproven assumptions
(\texttt{sorry}-free).  The statements quoted in this paper correspond to
the Lean theorems named in \S\ref{sec:mech}.  In the mathematical sections
below we therefore favour explaining the \emph{constructions} over spelling
out routine bookkeeping, which the formalization pins down exactly.

\paragraph{A caveat on independent evidence.}
During the development, bounded breadth-first searches over the state space
were used to hunt for counterexamples and to discover candidate move
sequences.  No counterexample to Theorem~\ref{thm:main-intro} was ever
found; all search results were subsequently superseded by machine-checked
proofs.  Claims below of the form ``$X$ is unreachable'' are proved
(Proposition~\ref{prop:sharp}); claims of the form ``the na\"ive strategy
fails'' are proved where stated and otherwise serve only as motivation.

\section{The game}\label{sec:model}

\begin{definition}[States and moves]\label{def:moves}
Fix $a,b,c\ge 1$ (the \emph{false sum}, written $\lie{a}{b}{c}$).  A
\emph{state} is a finite multiset of positive integers, called
\emph{balls}.  The legal moves are:
\begin{itemize}
\item \textbf{normal split}: replace one ball $n$ with $n\ge 2$ and
$n\ne c$ by the two balls $\lfloor n/2\rfloor,\ \lceil n/2\rceil$;
\item \textbf{false split}: replace one ball of value $c$ by the two balls
$a,\ b$;
\item \textbf{normal merge}: replace two balls $x,y$ with
$\{x,y\}\ne\{a,b\}$ (as multisets) by the ball $x+y$;
\item \textbf{false merge}: replace the two balls $a,\ b$ by the ball $c$.
\end{itemize}
We write $S\move T$ if $T$ is obtained from $S$ by one move, and $\reach$
for the reflexive--transitive closure.
\end{definition}

The player's false belief is \emph{forced}: the ball $c$ cannot be halved
(its only split is the false one), and the pair $\{a,b\}$ cannot merge to
$a+b$ (its only merge is the false one).  Note the model is manifestly
symmetric in $a\leftrightarrow b$, a fact we prove and use
(Lemma~\ref{lem:swap}).

\begin{definition}[Parameters]\label{def:params}
Set $\delta=a+b-c\in\Z$, \ $g=|\delta|$, \ $H=\max(a+b,c)$, and $M=H+1$.
\end{definition}

\begin{example}
Classic ($9+10=21$): $g=2$, $H=21$, $M=22$.  The false split lowers the
total by $2$; the false merge raises it by $2$.  For $6+6=3$: $g=9$,
$M=13$.  For $1+3=3$: $g=1$, $M=5$---so Theorem~\ref{thm:main-intro} asserts
that above $5$ \emph{every} pair is solvable.  Table~\ref{tab:examples}
lists more.
\end{example}

\begin{table}[t]
\centering
\begin{tabular}{lccl}
\toprule
false sum & $g$ & $M$ & character \\
\midrule
$9+10=21$ & $2$ & $22$ & Classic; deficient ($a+b<c$) \\
$3+3=5$ & $1$ & $7$ & excessive, both legs $<c$ \\
$8+9=4$ & $13$ & $18$ & one leg of the form $c\cdot 2^k$ \\
$15+15=4$ & $26$ & $31$ & large legs, not powers \\
$6+6=3$ & $9$ & $13$ & trap: $a=b=2c$ \\
$3+6=3$ & $6$ & $10$ & the pathological $\lie{c}{2c}{c}$ \\
$2+2=2$ & $2$ & $5$ & trap: $a=b=c$ \\
$4+6=2$ & $8$ & $11$ & $c=2$, both legs even \\
$1+3=3$ & $1$ & $5$ & doubly degenerate $\lie{1}{c}{c}$ \\
$2+3=1$ & $4$ & $6$ & $c=1$ \\
\bottomrule
\end{tabular}
\caption{Sample false sums with their inaccuracy $g$ and threshold $M$.
Each is solvable for all $s,t\ge M$ with $g\mid t-s$
(Theorem~\ref{thm:main-intro}).}\label{tab:examples}
\end{table}

\section{Necessity and sharpness}\label{sec:necessity}

\begin{proposition}[Necessity]\label{prop:necessity}
If $\ms{s}\reach\ms{t}$ then $g\mid t-s$.
\end{proposition}

\begin{proof}
Normal moves preserve the total of a state; the false split changes it by
$(a+b)-c=\delta$ and the false merge by $-\delta$.  Hence the total is
constant modulo $g$ along any play, and the totals of the two singleton
states are $s$ and $t$.
\end{proof}

The formalization proves this for an arbitrary finite \emph{set} of
simultaneous false sums, with $g$ the gcd of the individual inaccuracies;
only sufficiency is specific to a single lie.

\begin{proposition}[Sharpness at Classic]\label{prop:sharp}
For $\lie{9}{10}{21}$: \ $\ms{21}\not\reach\ms{23}$, although
$21,23\ge H=21$ and $2\mid 23-21$.
\end{proposition}

\begin{proof}
Call a state \emph{good} if all its balls are positive and either its total
is $19$ or it is exactly $\ms{21}$.  We claim goodness is preserved by every
move; since $\ms{23}$ is not good, this proves the proposition.  From
$\ms{21}$, the only available move is the false split ($21=c$ cannot halve;
nothing can merge), yielding $\ms{9,10}$ of total $19$: good.  From a good
state of total $19$: a false split would require a ball $21>19$, impossible;
a false merge requires balls $9$ and $10$, and since $9+10=19$ is the whole
total, the state is exactly $\ms{9,10}$ and the result is exactly
$\ms{21}$: good.  All other moves are normal and preserve the total $19$.
\end{proof}

Thus the guarantee of Theorem~\ref{thm:main-intro} fails at $H$ itself, and
$M=H+1$ is the optimal uniform threshold expressible in terms of $H$.  (For
particular false sums a smaller threshold may hold; we do not pursue the
per-configuration optimum here.)

\section{Sufficiency I: reduction to two pumps}\label{sec:pumps}

Everything below serves the converse direction.  The first step is a
routine reduction.

\begin{lemma}[Pump reduction]\label{lem:pumps}
Suppose the two \emph{pumps} hold for a given $\lie{a}{b}{c}$:
\[
\text{(climb)}\quad \forall n\ge M:\ \ms{n}\reach\ms{n+g},
\qquad
\text{(descend)}\quad \forall n\ge M:\ \ms{n+g}\reach\ms{n}.
\]
Then $\ms{s}\reach\ms{t}$ for all $s,t\ge M$ with $g\mid t-s$.
\end{lemma}

\begin{proof}
Iterate the appropriate pump $|t-s|/g$ times; every intermediate value stays
$\ge M$.
\end{proof}

Note that reachability is \emph{not} symmetric move-by-move (a normal merge
$\{3,7\}\to 10$ reverses only via a longer route), so climb and descend
genuinely require separate constructions.  The remainder of the paper
discharges the two pumps for every $(a,b,c)$, by cases.

\section{Sufficiency II: the all-ones hub}\label{sec:hub}

Throughout this section, ``units'' are balls of value $1$.  The hub
strategy for $\ms{n}\reach\ms{n\pm g}$ is:
\begin{enumerate}
\item \emph{scatter} $\ms{n}$ into the pile of units $1^{n}$;
\item change the pile size by $\pm g$ using one false move
(\emph{gain}: create a ball $c$, false-split it into $\{a,b\}$, re-scatter,
for $a+b>c$; \emph{lose}: assemble $a$ and $b$, false-merge to $c$,
re-scatter; roles swap when $a+b<c$);
\item \emph{rebuild} the single ball from the pile.
\end{enumerate}
Each step must dodge the two special values: a ball $c$ is \emph{locked}
(it cannot halve), and the pair $\{a,b\}$ is \emph{poisoned} (it merges only
falsely).  The precise dodges depend on the parameters.

\subsection{The deficient case \texorpdfstring{$a+b<c$}{a+b<c}}

Here $a,b<c$, and every value below $c$ is \emph{free}: its halves stay
below $c$, so it scatters to units by repeated normal splits, and units
re-assemble into any value $k\le\max(a,b)$ (merging a unit onto $j<k$
can only form the poisoned pair if $\{1,j\}=\{a,b\}$, which forces
$\max(a,b)=j<k\le \max(a,b)$---absurd).  The climb pump on the window
$n\in[c+1,\,2c-2]$ is then: split $n$ into two sub-$c$ halves, scatter both,
assemble $a$ and $b$, false-merge to $c$ ($+g$), and reel the remaining
units onto the $c$ (safe: $c>\max(a,b)$, so $\{x,1\}$ with $x \geq c$ is
never poisoned).  The three boundary values $n\in\{2c-1,2c,2c+1\}$, whose
halves collide with the locked $c$, are handled by short explicit routes,
and a halving recursion (split, pump the half, re-merge) reduces every
$n\ge M$ to the window.  The descend pump is dual (false-split a freshly
built $c$).

Two edges of this case need extra care, both caused by unit legs.  If
$a=1$ the pair $\{1,b\}$ is poisoned, so re-assembly must \emph{skip} the
value $b+1$: build $b$ and a separate $2$, merge them to $b+2$ (legal:
$\{b,2\}$ is never $\{1,b\}$), and continue.  If $a=b=1$, units cannot merge
\emph{at all} ($\{1,1\}$ is the false pair, merging to $c$!); the
construction instead maintains \emph{2-seeds}---balls of value $2$ produced
by splitting odd values---and grows from those.

\subsection{The excessive case with small legs}

For $a+b>c$ and both legs $<c$ the same hub runs with the roles of the
false moves exchanged (the false \emph{split} now gains $g$).  This case
includes every ``cluster structure'' of the boundary values and is fully
generic; the formal statement covers all $2\le a,b<c<a+b$.

When exactly one leg is small ($2\le a<c\le b$, say), scattering the big leg
may jam on locked $c$'s, and the following device---used pervasively from
here on---resolves it.

\begin{lemma}[Reservoir bumping]\label{lem:bump}
Let $c\ge 3$ and suppose $\{1,c\}\ne\{a,b\}$.  From any state containing a
ball $v\ge 1$ and $K\ge 1$ units, one can reach the state with $v+K$
units.
\end{lemma}

\begin{proof}[Proof sketch]
Halve $v$ repeatedly.  Whenever a locked $c$ appears, \emph{bump} it: merge
one unit onto it (legal by hypothesis), obtaining $c+1$, which is unlocked
and halves into two values $\le\lceil (c+1)/2\rceil\le c-1$, both free to
continue.  The unit invested is recovered because all fragments eventually
reach $1$.  A strong induction on the multiset of values makes this
terminate.
\end{proof}

Thus a single unit ``unlocks'' everything---the problem is only ever
\emph{obtaining the first units}, which is exactly what the trap families
prevent.

\section{Scatterability, and the trap dichotomy}\label{sec:scat}

For $a+b>c$ with both legs $\ge c$ we need to know which single balls can
produce units on their own.

\begin{definition}[Scatterable values]\label{def:scat}
For $c\ge 3$, let $\Scat_c\subseteq\N$ be the least set such that:
$v\in\Scat_c$ whenever $c<v<2c$; and $v\in\Scat_c$ whenever $v\ge 2c$ and
($\lfloor v/2\rfloor\in\Scat_c$ or $\lceil v/2\rceil\in\Scat_c$).
\end{definition}

A value $c<v<2c$ splits into a half $<c$ (which scatters freely into a
reservoir of units) plus a remainder which is then bumped
(Lemma~\ref{lem:bump}); a larger scatterable value halves towards the
window.  Hence, under the safety hypothesis of Lemma~\ref{lem:bump}, every
$v\in\Scat_c$ satisfies $\ms{v}\reach 1^{v}$ on its own.  The values
\emph{not} in $\Scat_c$ are precisely characterized:

\begin{lemma}[Dichotomy]\label{lem:dichotomy}
Let $c\ge 3$.  Every $v\ge c$ satisfies: $v=2^k c$ for some $k\ge 0$, or
$v\in\Scat_c$.
\end{lemma}

\begin{proof}
Strong induction on $v$.  If $c\le v<2c$: either $v=c=2^0c$, or $v$ is in
the window.  If $v\ge 2c$ is even, its two halves equal $v/2\ge c$; by
induction $v/2$ is a power multiple (making $v$ one) or scatterable (making
$v$ scatterable).  If $v\ge 2c$ is odd, its halves are consecutive integers
$f,f+1\ge c$; they cannot \emph{both} be of the form $2^ic$, since
$2^jc-2^ic=1$ would force $c\mid 1$.  Apply induction to a half that is not
a power multiple: it is scatterable, hence so is $v$.
\end{proof}

Consequently, if at least one leg is not of the form $2^kc$, that leg
scatters on its own; a false split $c\to\{a,b\}$ then converts any locked
$c$ into a scatterable pile, every value scatters, and the hub closes the
configuration.  (This subsumes, in one stroke, all ``bands'' of leg sizes.)
What remains is exactly:

\begin{center}
\emph{both legs of the form $2^kc$}\qquad and\qquad
\emph{the small targets $c\in\{1,2\}$ and unit legs.}
\end{center}

\section{Sufficiency III: the traps}\label{sec:traps}

Let now $a=2^pc$ and $b=2^qc$, $c\ge 3$.  These are genuine traps for the
hub: for instance in $\lie{2}{2}{2}$ a lone ball $2$ generates only
multisets of $2$'s (its only split is $2\to\{2,2\}$ and two $2$'s merge
only falsely back to $2$), and the analogous $c$-multiple phenomena occur
for all these families.  We do not need any unreachability statement,
however; we simply replace the commodity: \emph{hub on copies of $c$
instead of units}.

\begin{lemma}[Peeling]\label{lem:peel}
Let $a=2^pc$, $b=2^qc$, $c\ge 2$, and suppose $\{a,b\}\ne\{c,2c\}$.  Then
for every $v\ge c+1$:\ \ $\ms{v}\reach\ms{c,\,v-c}$.
\end{lemma}

\begin{proof}[Proof sketch]
Strong induction.  For $v\le 2c-2$ both halves of $v$ are $<c$, hence free:
scatter everything to units and re-assemble $c$ and $v-c$ (all merges among
sub-$c$ values and units are unpoisoned, since $a,b\ge c$).  The boundary
values $2c-1,2c,2c+1$ split directly into $\{c-1,c\},\{c,c\},\{c,c+1\}$.
For $v\ge 2c+2$: split $v$, peel a $c$ off the lower half by induction, and
re-merge the two leftovers $\lfloor v/2\rfloor-c$ and $\lceil v/2\rceil$.
This merge is poisoned only if
$\{\lfloor v/2\rfloor-c,\lceil v/2\rceil\}=\{a,b\}$, which forces
$|a-b|\in\{c,c+1\}$.  For legs $2^pc\ne 2^qc$ the gap $|a-b|$ is a nonzero
multiple of $c$, so only $|a-b|=c$ is possible, and $2^q-2^p=1$ then forces
$(p,q)=(0,1)$ (Lemma~\ref{lem:powgap}): that is, $\{a,b\}=\{c,2c\}$,
excluded by hypothesis.
\end{proof}

Given peeling, the pumps are short.  \emph{Climb}: peel one $c$ off
$\ms{n}$, false-split it into $\{a,b\}$ ($+g$), and merge everything back
(both merges avoid the poisoned pair because the accumulating ball exceeds
$b$).  \emph{Descend}: peel $2^p+2^q$ copies of $c$ off $\ms{n+g}$, rebuild
the two legs from the copies, false-merge $\{a,b\}\to c$ ($-g$), and merge
back.  Rebuilding a leg $2^kc$ from $2^k$ copies of $c$ uses \emph{binary
doubling}, $\ms{x,x}\to\ms{2x}$, which can never hit the poisoned pair when
$a\ne b$; when $a=b$ (so $p=q\ge 1$) one instead merges copies
sequentially, $\ms{c,jc}\to\ms{(j+1)c}$, unpoisoned since $jc<a$ and
$c\neq a$.  The diagonal $a=b=c$ is simplest of all: climb peels one $c$
and false-splits it; descend peels two and false-merges them.

\subsection{The pathological shape \texorpdfstring{$\lie{c}{2c}{c}$}{(c,2c,c)}}
\label{sec:peel4c}

One shape survives the analysis above: $a=c$, $b=2c$ (``$c+2c=c$'', e.g.\
$3+6=3$), where the peeling recursion of Lemma~\ref{lem:peel} jams, at
exactly one ball value.  At $v=4c$ the split gives $\{2c,2c\}$, peeling a
$c$ off one half leaves $\{c,2c\}$---precisely the poisoned pair---and no
local reordering avoids it.  The escape is an eleven-move excursion that
climbs to total $8c$, forges the otherwise unobtainable ball $3c$ by
splitting a $6c$, and rides two false merges back down:
\[
\begin{array}{r@{\;}c@{\;}l@{\qquad}ll}
\ms{4c} &\move& \ms{2c,2c} & \text{split $4c$} & \text{(total $4c$)}\\
        &\move& \ms{c,c,2c} & \text{split $2c$} \\
        &\move& \ms{c,c,2c,2c} & \text{false split } c\to\{c,2c\} & (+g) \\
        &\move& \ms{c,c,2c,2c,2c} & \text{false split } c\to\{c,2c\} & (+g)\ \text{(total $8c$)} \\
        &\move& \ms{c,c,2c,4c} & \text{merge }\{2c,2c\} \\
        &\move& \ms{c,c,6c} & \text{merge }\{2c,4c\} \\
        &\move& \ms{c,c,3c,3c} & \text{split } 6c \\
        &\move& \ms{c,3c,4c} & \text{merge }\{c,3c\} \\
        &\move& \ms{c,2c,2c,3c} & \text{split } 4c \\
        &\move& \ms{c,2c,3c} & \text{false merge }\{c,2c\}\to c & (-g) \\
        &\move& \ms{c,3c} & \text{false merge }\{c,2c\}\to c & (-g)\ \text{(total $4c$)}
\end{array}
\]
With this single gadget grafted in at $v=4c$, the peeling recursion---and
with it both pumps---goes through for $\lie{c}{2c}{c}$, for every $c\ge 2$.

\subsection{The doubly degenerate family \texorpdfstring{$\lie{1}{c}{c}$}{(1,c,c)}}
\label{sec:onecc}

The false sum $1+c=c$ (any $c\ge2$; $g=1$) is a trap of a different kind:
$a=1$ makes the pair $\{1,c\}$ poisoned, so the ubiquitous ``merge a unit
onto a locked $c$'' bump of Lemma~\ref{lem:bump} is illegal, and the ball
$c$ regenerates under its own false split ($c\to\{1,c\}$).  The climb pump
is still easy (peel a $c$---Lemma~\ref{lem:peel} applies, its recursion
never jams here---false-split it, and re-merge, for a net $+1$).  The
descend pump must \emph{lose exactly one} from $\ms{n+1}$, and the
construction is a recursion we call \textsf{loseOne}: peel copies of $c$
off the ball until the remainder drops below $c$, scatter the (free)
remainder into units, ride \emph{one} false merge $\{1,c\}\to c$
(erasing a single unit, $-1$), and re-assemble---the re-assembly dodging
$\{1,c\}$ throughout because the accumulated $c$-copies merge with each
other and with sub-$c$ values only.  One divisibility wrinkle remains: if
$c\mid n+1$ the peeling bottoms out with no remainder at all (an all-$c$'s
pile).  The fix is to \emph{overshoot}: climb once more to $n+2\equiv
1\pmod c$, peel down to the remainder $c+1$, scatter it, and ride
\emph{two} false merges.  For $c=2$ this last route needs units from a pile
of $2$'s, obtained not by false-splitting (which regenerates) but by the
\emph{odd-ball core}: merge three $2$'s into $6$, split $6\to\{3,3\}$, and
split each $3\to\{1,2\}$---manufacturing two units by \emph{normal} moves
alone.

\section{Sufficiency IV: the small targets \texorpdfstring{$c=1$ and $c=2$}{c=1 and c=2}}
\label{sec:smallc}

\paragraph{$c=1$.}
Nothing is ever locked: normal splitting requires $n\ge2\ne c$, so every
ball halves freely, and only the poisoned pair $\{a,b\}$ needs dodging.
The hub runs essentially unobstructed; the one point of care is that for
$a=1$ the re-assembly must skip the poisoned $\{1,b\}$ exactly as in the
deficient case (build $b-1$ and a spare $2$, or two $2$'s when $b=2$).

\paragraph{$c=2$.}
Here units are a scarce resource: the ball $2$ is locked (it false-splits
into $\{a,b\}$), so an even ball can decompose into $2$'s but no further.
The construction hubs on \emph{copies of $2$}, with a two-route peeling
lemma ($\ms{v}\reach\ms{2,v-2}$ for $v\ge3$: when the standard recursion
would produce the poisoned pair, peel off the \emph{other} half, whose
residue differs by at most one and is therefore safe), and rebuilds legs
from $2$'s.  Odd legs are manufactured by \emph{splitting} an even ball
($\ms{2,2,2}\to\ms{6}\to\ms{3,3}\to\ms{1,2,3}$, a gadget that also yields a
spare unit and a spare $2$, both of which are absorbed harmlessly).  The
analysis splits into the four parity classes of $(a,b)$, plus the
degenerate configurations where a leg equals $2$ (then $a=c$, and every
$\{x,2\}$ with $x=b$ is poisoned) or equals $1$; each is closed by a
variant of the same two devices.  We spare the reader the case ledger---the
formalization contains ten theorems whose union is: \emph{every
$\lie{a}{b}{2}$ with $a+b>2$ satisfies both pumps.}

\section{The assembly}\label{sec:assembly}

Three glue facts combine the families into a single statement.

\begin{lemma}[Swap]\label{lem:swap}
Reachability, $g$, and $M$ are invariant under $\lie{a}{b}{c}\mapsto
\lie{b}{a}{c}$.
\end{lemma}

\begin{lemma}[Power gap]\label{lem:powgap}
If $2^q=2^p+1$ then $p=0$ and $q=1$.  Consequently, for legs $a=2^pc\ne
b=2^qc$: \ $b=a+c$ forces $\{a,b\}=\{c,2c\}$, and $b=a+c+1$ is impossible
for $c\ge2$ (it would give $c\mid 1$).
\end{lemma}

\begin{theorem}[Characterization; formally verified]\label{thm:main}
For all $a,b,c\ge1$ and all $s,t\ge M$:
\[
\ms{s}\reach\ms{t}\iff g\mid t-s.
\]
\end{theorem}

\begin{proof}[Proof (assembly)]
Necessity is Proposition~\ref{prop:necessity}.  For sufficiency, by
Lemma~\ref{lem:swap} assume $a\le b$ and dispatch:
if $a+b=c$, then $g=0$ and $t=s$.  If $a+b<c$: the deficient hub
(\S\ref{sec:hub}), in its three variants ($2\le a$; $a=1<b$; $a=b=1$).  If
$a+b>c$: for $c=1$ and $c=2$, \S\ref{sec:smallc}; for $c\ge3$ with $a=1$,
either $b>c$ (unit-leg hub, a bumping variant) or $b=c$
(\S\ref{sec:onecc}); for $2\le a,b<c$, the excessive hub; for $a<c\le b$,
the small-leg bootstrap; and for $c\le a\le b$, apply
Lemma~\ref{lem:dichotomy} to both legs: if either is scatterable, the hub
closes it (\S\ref{sec:scat}); otherwise both are powers $2^pc,2^qc$ and the
trap constructions apply---the diagonal $p=q=0$, the equal case $p=q\ge1$,
the shape $\{c,2c\}$ (via \S\ref{sec:peel4c}, whose applicability is
exactly delimited by Lemma~\ref{lem:powgap}), and the generic
$p\ne q$ (Lemma~\ref{lem:peel}).  Every branch discharges both pumps, and
Lemma~\ref{lem:pumps} concludes.
\end{proof}

\section{The formalization}\label{sec:mech}

The entire development is verified in Lean~4 (v4.31.0) \cite{lean4} using
only the core library---no \textsf{Mathlib}---in a single file
(\texttt{lean/YaStupid.lean}, $7{,}462$ lines, $328$ theorems) in the
repository accompanying this paper.\footnote{\url{https://github.com/emblaisdell/yastupid}}
States are lists of naturals quotiented by permutation at the level of the
step relation; the four moves of Definition~\ref{def:moves} are an
inductive family \texttt{Local}, one step applies a \texttt{Local} move to
a sub-multiset (\texttt{Step}), and \texttt{Reach} is its
reflexive--transitive closure.  The headline statements are:

\begin{center}
\small
\begin{tabular}{ll}
\toprule
paper & Lean theorem \\
\midrule
Theorem~\ref{thm:main} ($\Leftarrow$, all $a,b,c$) & \texttt{single\_sufficiency\_all} \\
Theorem~\ref{thm:main} (iff form) & \texttt{single\_characterization} \\
Proposition~\ref{prop:necessity} (any \# of lies) & \texttt{reach\_congr} \\
Proposition~\ref{prop:sharp} & \texttt{classic\_trap} \\
Lemma~\ref{lem:pumps} & \texttt{sufficiency\_of\_pumps} \\
Lemma~\ref{lem:dichotomy} & \texttt{pow\_or\_scat} \\
Lemma~\ref{lem:peel} & \texttt{peelcG} \\
\S\ref{sec:peel4c} gadget & \texttt{peel4c} \\
\S\ref{sec:onecc} recursion & \texttt{loseOne} \\
Lemma~\ref{lem:swap} & \texttt{reach\_swap}, \texttt{suff\_swap} \\
Lemma~\ref{lem:powgap} & \texttt{pow2\_eq\_succ} \\
\bottomrule
\end{tabular}
\end{center}

The file compiles in about $40$ seconds and every theorem is checked (via
\texttt{\#print axioms}) to depend only on Lean's three standard classical
axioms \texttt{propext}, \texttt{Classical.choice}, \texttt{Quot.sound}; in
particular the development is free of \texttt{sorry}.  Beyond the master
theorems, the file contains machine-checked closed-form corollaries for the
concrete false sums of Table~\ref{tab:examples} (e.g.\
\texttt{solvable\_6\_6\_3}: all $s,t\ge13$ with $9\mid t-s$ are
interreachable under $6+6=3$).

Formalization was not an afterthought: several ``obvious'' claims failed
under mechanization and reshaped the mathematics.  The all-ones hub was
first believed to work whenever the pumps could bridge $g$-gaps; the trap
families refuted this.  The peeling recursion was first stated without the
$\{c,2c\}$ exclusion; Lean's unprovable goal at the re-merge step is what
isolated the pathological shape of \S\ref{sec:peel4c} and the exact
arithmetic of Lemma~\ref{lem:powgap}.

\section{Concluding remarks}\label{sec:conclusion}

\begin{itemize}
\item \textbf{Multiple lies.}  For a finite set of simultaneous false sums
the congruence obstruction (with $g$ the gcd of the inaccuracies) is proved
in the same formalization, and we conjecture the analogue of
Theorem~\ref{thm:main} with $H$ the maximum of the individual $H$'s.
The interaction of several locked values and poisoned pairs makes the
trap taxonomy considerably richer, and we leave it open.
\item \textbf{Per-configuration thresholds.}  $M=H+1$ is optimal as a
uniform function of $H$ (Proposition~\ref{prop:sharp}), but individual
false sums can admit smaller thresholds; characterizing the exact minimum
for each $(a,b,c)$ is open.
\item \textbf{Beyond singletons.}  Our characterization concerns
single-ball endpoints.  General multiset-to-multiset reachability adds a
second obstruction (the number of balls interacts with the totals) and
appears substantially harder.
\end{itemize}

\subsection*{Acknowledgements}
The move sequences underlying several constructions were first located by
bounded breadth-first search before being generalized and verified; we
thank the machine for its patience.

\begin{thebibliography}{1}
\bibitem{lean4}
L.~de~Moura and S.~Ullrich.
\newblock The {Lean} 4 theorem prover and programming language.
\newblock In \emph{Automated Deduction -- CADE 28}, LNCS 12699, pp.~625--635.
Springer, 2021.
\end{thebibliography}

\end{document}
